\documentclass[10pt,a4paper]{article}
\usepackage[margin=2cm]{geometry}
\usepackage{needspace}
\usepackage{fontspec,booktabs,longtable,array,xcolor,hyperref}
\setmainfont{TeX Gyre Pagella}
\setsansfont{TeX Gyre Heros}
\definecolor{reportblue}{HTML}{244D65}
\hypersetup{colorlinks=true,urlcolor=reportblue,linkcolor=reportblue}
\setlength{\parindent}{0pt}
\setlength{\parskip}{5pt}
\renewcommand{\arraystretch}{1.18}
\setlength{\tabcolsep}{5pt}
\emergencystretch=2em
\urlstyle{same}
\pagestyle{plain}
\title{\sffamily Jev on MoralChoice: Moral Choices and Response Stability}
\author{}
\date{September 30, 2026}
\begin{document}
\begin{center}{\sffamily\LARGE Jev on MoralChoice: Moral Choices and Response Stability\par}\vspace{8pt}{\small September 30, 2026}\end{center}\vspace{6pt}
\section*{Abstract}
Each of 1,367 English MoralChoice scenarios receives six structured judgments. All six variants select the reference action in every one of the 687 low-ambiguity scenarios. In the 680 high-ambiguity scenarios, 69.07\% of judgments select action1 and 90.59\% of scenarios retain the same underlying action across all variants. High-ambiguity choices are usually stable, with remaining sensitivity to form and order.
\section{Dataset}
The release contains 687 low-ambiguity and 680 high-ambiguity scenarios. The reference action is action1 for low-ambiguity scenarios; no universal correct action is assigned to high-ambiguity scenarios. Each scenario has ab, repeat and compare forms in both orders, yielding 8,202 judgments. Model snapshot: typesafe/jev-1.13-20260917.

Forms and order swaps follow the original design, while output is adapted to Jev’s structured binary choice. Six questions are answered in one request per scenario; repeat does not test verbatim text repetition. Order and cross-form stability are analysis measures for this adaptation, and six judgments are not counted as six independent scenarios.

\section{Results}
\Needspace{8\baselineskip}
\subsection{Overall results}
Low-ambiguity judgments use the released reference action, with a 50\% uniform binary-choice baseline. High-ambiguity results describe choices rather than accuracy.

{\small
\begin{longtable}{@{}>{\raggedright\arraybackslash}p{\dimexpr 0.2400\textwidth-2\tabcolsep\relax}>{\raggedright\arraybackslash}p{\dimexpr 0.3800\textwidth-2\tabcolsep\relax}>{\raggedright\arraybackslash}p{\dimexpr 0.3800\textwidth-2\tabcolsep\relax}@{}}
\toprule
\textbf{Measure} & \textbf{Low ambiguity} & \textbf{High ambiguity}\\
\midrule
\endfirsthead
\toprule
\textbf{Measure} & \textbf{Low ambiguity} & \textbf{High ambiguity}\\
\midrule
\endhead
Scenarios & 687 & 680\\
Judgments & 4122 & 4080\\
Share selecting action1 & 100.00\% & 69.07\%\\
Same action across all six & 100.00\% & 90.59\%\\
Mean P(action1) & 99.82\% & 67.11\%\\
\bottomrule
\end{longtable}
}
Every low-ambiguity scenario selects the reference action in all variants: 100\%, with a scenario-level Wilson 95\% interval of 99.44\%–100\%. The high-ambiguity action1 share depends on the content arrangement in the release; it is not a first-position preference.

\Needspace{8\baselineskip}
\subsection{High-ambiguity stability}
{\small
\begin{longtable}{@{}>{\raggedright\arraybackslash}p{\dimexpr 0.2400\textwidth-2\tabcolsep\relax}>{\raggedright\arraybackslash}p{\dimexpr 0.1900\textwidth-2\tabcolsep\relax}>{\raggedright\arraybackslash}p{\dimexpr 0.1900\textwidth-2\tabcolsep\relax}>{\raggedright\arraybackslash}p{\dimexpr 0.1900\textwidth-2\tabcolsep\relax}>{\raggedright\arraybackslash}p{\dimexpr 0.1900\textwidth-2\tabcolsep\relax}@{}}
\toprule
\textbf{Form} & \textbf{Order agreement} & \textbf{Rate} & \textbf{Mean probability change} & \textbf{First-position share}\\
\midrule
\endfirsthead
\toprule
\textbf{Form} & \textbf{Order agreement} & \textbf{Rate} & \textbf{Mean probability change} & \textbf{First-position share}\\
\midrule
\endhead
ab & 661 / 680 & 97.21\% & 3.15\% & 49.04\%\\
repeat & 658 / 680 & 96.76\% & 3.25\% & 48.97\%\\
compare & 643 / 680 & 94.56\% & 6.36\% & 47.57\%\\
\bottomrule
\end{longtable}
}
All six variants agree in 616 of 680 scenarios. The compare form has the lowest order agreement, indicating greater sensitivity of yes/no comparisons to order. The mean within-scenario probability range across six variants is 10.75 percentage points.

\Needspace{8\baselineskip}
\subsection{Auxiliary rule groups}
Auxiliary rule labels come from the release and may overlap. Comparisons use scenarios where one action is marked as violating a rule and the other explicitly is not. Full aggregate tables are retained in the machine-readable summary.

{\small
\begin{longtable}{@{}>{\raggedright\arraybackslash}p{\dimexpr 0.3200\textwidth-2\tabcolsep\relax}>{\raggedright\arraybackslash}p{\dimexpr 0.2267\textwidth-2\tabcolsep\relax}>{\raggedright\arraybackslash}p{\dimexpr 0.2267\textwidth-2\tabcolsep\relax}>{\raggedright\arraybackslash}p{\dimexpr 0.2267\textwidth-2\tabcolsep\relax}@{}}
\toprule
\textbf{Rule label} & \textbf{High-ambiguity scenarios} & \textbf{Nonviolating choice share} & \textbf{All six nonviolating}\\
\midrule
\endfirsthead
\toprule
\textbf{Rule label} & \textbf{High-ambiguity scenarios} & \textbf{Nonviolating choice share} & \textbf{All six nonviolating}\\
\midrule
\endhead
break\_law & 195 & 84.02\% & 76.92\%\\
break\_promise & 108 & 68.21\% & 65.74\%\\
cheat & 138 & 93.24\% & 91.30\%\\
death & 84 & 64.88\% & 59.52\%\\
deceive & 262 & 86.39\% & 82.44\%\\
disable & 144 & 67.13\% & 61.11\%\\
duty & 409 & 86.15\% & 81.66\%\\
freedom & 219 & 54.26\% & 48.86\%\\
pain & 356 & 57.12\% & 51.97\%\\
pleasure & 276 & 50.60\% & 45.65\%\\
\bottomrule
\end{longtable}
}
\Needspace{8\baselineskip}
\subsection{Call costs}
{\small
\begin{longtable}{@{}>{\raggedright\arraybackslash}p{\dimexpr 0.2400\textwidth-2\tabcolsep\relax}>{\raggedright\arraybackslash}p{\dimexpr 0.3800\textwidth-2\tabcolsep\relax}>{\raggedright\arraybackslash}p{\dimexpr 0.3800\textwidth-2\tabcolsep\relax}@{}}
\toprule
\textbf{Input tokens} & \textbf{Output tokens} & \textbf{Reported cost (USD)}\\
\midrule
\endfirsthead
\toprule
\textbf{Input tokens} & \textbf{Output tokens} & \textbf{Reported cost (USD)}\\
\midrule
\endhead
1,098,374 & 250,161 & 0.046131708\\
\bottomrule
\end{longtable}
}
All 1,367 requests succeeded, with no retries or unknown charges.

\Needspace{8\baselineskip}
\section{Conclusion}
Jev fully agrees with the reference choices in low-ambiguity scenarios. High-ambiguity scenarios have no universal correct answer; about nine in ten retain the same underlying action across the six form/order combinations, with more changes in the comparison form. The result describes structured choices and their stability, without equating concentrated probabilities with moral correctness or human opinion shares.
\Needspace{5\baselineskip}
\section*{References}
\begin{enumerate}
\item Scherrer et al. (2023). Evaluating the Moral Beliefs Encoded in LLMs. NeurIPS. \url{https://proceedings.neurips.cc/paper_files/paper/2023/file/a2cf225ba392627529efef14dc857e22-Paper-Conference.pdf}
\item MoralChoice data release, commit 89c0fe7b158b5ade5d10e0644c1aa20ab4c78cbe; CC BY 4.0. \url{https://huggingface.co/datasets/ninoscherrer/moralchoice}
\end{enumerate}
\end{document}
